% Fichier généré automatiquement : explication_detaill_e.tex
% Document unifié et nettoyé : explication_detaill_e.tex
\documentclass[a4paper,11pt]{article}
\usepackage[utf8]{inputenc}
\usepackage[T1]{fontenc}
\usepackage[french]{babel}
\usepackage{graphicx}
\usepackage{amsmath,amssymb}
\usepackage{geometry}
\geometry{margin=2.5cm}
\usepackage{hyperref}
\usepackage{caption}

\title{Explication détaillée — Logique floue et ensembles flous}
\author{Généré par l'assistant}
\date{\today}

\begin{document}
\maketitle

\section*{Objectif}
Ce document unifié rassemble toutes les explications demandées durant la session : définitions d'ensembles flous, fonctions d'appartenance (triangulaire et trapézoïdale), interprétations pédagogiques (versions pour un(e) ado de 15 ans et pour un enfant de 8 ans), exemples numériques, éléments pratiques basés sur MovieLens, et les illustrations produites. Le document ne contient pas de code ; les scripts utilisés pour générer les figures existent séparément dans le dépôt.

\section{Rappels simples}
\subsection{Univers et ensembles}
L'\emph{univers} $X$ est l'ensemble des valeurs possibles (par exemple toutes les notes : $[0,5]$ ou celles présentes dans MovieLens). Un \emph{ensemble flou} $A$ est un concept (par exemple «bon film») défini par une fonction d'appartenance $\mu_A$ qui associe à chaque $x\in X$ un degré $\mu_A(x)\in[0,1]$.

\subsection{Fonction d'appartenance}
La fonction $\mu_A:X\to[0,1]$ donne pour chaque valeur $x$ un nombre entre 0 et 1 : 0 = pas du tout, 1 = complètement. On peut lire $\mu_A(x)$ comme «à quel point $x$ appartient au concept $A$».

\section{Explications pédagogiques}
\subsection*{Pour un(e) ado de 15 ans}
Imagine un curseur entre 0 et 1 pour chaque valeur (par ex. une note) qui indique dans quelle mesure la valeur correspond au concept ("bon", "mauvais", ...). Les fonctions triangulaires et trapézoïdales sont des formes simples pour définir ces curseurs.

\subsection*{Pour un enfant de 8 ans}
Pense à une lampe : 0 = éteinte, 1 = allumée à fond. Dire qu'un film est "un peu bon" revient à laisser la lampe à moitié allumée. Un triangle correspond à une lampe qui s'allume au maximum en un seul point (un pic). Un trapèze, c'est une lampe qui reste complètement allumée sur un petit intervalle (un plateau).

\section{Fonctions d'appartenance : triangle et trapèze}
\subsection*{Triangulaire}
Une fonction triangulaire est définie par trois paramètres $(a,b,c)$, avec $a<b<c$ :
$$
\mu(x)=\begin{cases}
0, & x\le a,\\
\dfrac{x-a}{b-a}, & a<x\le b,\\
\dfrac{c-x}{c-b}, & b<x<c,\\
0, & x\ge c.
\end{cases}
$$
Le point $b$ est le sommet où $\mu(b)=1$. Le triangle est adapté quand on veut modéliser un seul point central d'appartenance maximale.

\subsection*{Trapézoïde}
Un trapèze est défini par quatre paramètres $(a,b,c,d)$, $a\le b\le c\le d$ :
$$
\mu(x)=\begin{cases}
0, & x\le a,\\
\dfrac{x-a}{b-a}, & a<x\le b,\\
1, & b<x\le c,\\
\dfrac{d-x}{d-c}, & c<x<d,\\
0, & x\ge d.
\end{cases}
$$
Le plateau $[b,c]$ contient les valeurs considérées comme entièrement appartenantes (tolérance / robustesse).

\section{Exemples numériques}
\subsection*{Interprétation simple}
Si l'on prend un triangle $(a,b,c)=(0.5,3.0,3.5)$ alors :
\begin{itemize}
    \item pour $x=2.0$ on a $\mu(2.0)=(2.0-0.5)/(3.0-0.5)=1.5/2.5=0.60$ ; la note 2.0 appartient à 60\% au concept défini par ce triangle.
    \item pour $x=3.0$ (le sommet) on a $\mu(3.0)=1.00$.
\end{itemize}
Pour un trapèze $(a,b,c,d)=(0.5,0.5,2.0,3.0)$, le plateau $[b,c]=[0.5,2.0]$ contient $x=2.0$ et donc $\mu(2.0)=1.0$.

\section{Illustrations}
Les figures générées pendant la session illustrent ces notions :
\begin{figure}[ht]
    \centering
    \includegraphics[width=0.9\linewidth]{figures/membership_examples.png}
    \caption{Comparaison de deux fonctions triangulaires et d'un trapézoïde (exemples simples). La verticale marque $x=3.0$ et les étiquettes indiquent les $\mu$ correspondants.}
\end{figure}

\begin{figure}[ht]
    \centering
    \includegraphics[width=0.95\linewidth]{figures/membership_moviLens_quantiles.png}
    \caption{Histogramme des notes MovieLens avec fonctions d'appartenance construites à partir de quantiles. Triangles = partition intuitive, trapèzes = intervalles de tolérance.}
\end{figure}

\section{Interprétation des figures (raccourci)}
\begin{itemize}
    \item Triangle $(2,3,5)$ : $\mu(3)=1$, $\mu(2.5)=0.5$.
    \item Triangle $(1,3.5,5)$ : sommet à 3.5 ; $\mu(3.0)\approx 0.80$.
    \item Trapèze $(1,2,4,5)$ : plateau $[2,4]$ avec $\mu=1$, donc $\mu(3)=1$.
    \item Sur MovieLens, les quantiles observés (exemple) étaient : min=0.50, q10=2.00, q25=3.00, q50=3.50, q75=4.00, q90=5.00. On peut construire des triangles et trapèzes à partir de ces repères pour obtenir des termes low/mid/high interprétables.
\end{itemize}

\section{Paramètres pratiques (extrait de la configuration)}
Exemple d'utilisation pratique (extrait de \texttt{config/fuzzy_config.yaml}) :
\begin{itemize}
    \item \textbf{genre\_preference} (univers $[0.0,1.0]$) :
        \begin{itemize}
            \item \emph{faible} : trapézoïdal paramètres [0.0, 0.0, 0.2, 0.45]
            \item \emph{moyenne} : triangulaire paramètres [0.3, 0.5, 0.7]
            \item \emph{forte} : trapézoïdal paramètres [0.55, 0.8, 1.0, 1.0]
        \end{itemize}
    \item \textbf{average\_rating} (univers $[0.5,5.0]$) : termes \emph{mauvaise}, \emph{correcte}, \emph{bonne}, \emph{excellente} avec paramètres choisis autour des quartiles/valeurs observées.
    \item \textbf{recommendation\_score} (sortie, univers $[0,1]$) : terme final réparti en plusieurs trapèzes/triangles pour produire une sortie interprétable.
\end{itemize}

\section{Questions fréquentes}
\begin{itemize}
    \item "La somme des $\mu$ vaut-elle 1 ?" — Non : les degrés d'appartenance sont indépendants et ne constituent pas une distribution de probabilité obligatoire.
    \item "Pourquoi choisir plateau (trapèze) ou pic (triangle) ?" — Trapèze = robustesse / tolérance ; triangle = sensibilité / précision.
    \item "Comment calibrer automatiquement ?" — Heuristiques : placer paramètres sur des percentiles (q25,q50,q75), ou optimiser par apprentissage pour améliorer les métriques de recommandation.
\end{itemize}

\section*{Conclusion et procédure de mise à jour}
Le présent fichier `explication\_detaill\_e.tex` est désormais le document unique d'explication. À chaque nouvelle demande d'explication, je l'intègrerai dans ce même fichier et pourrai le recompiler en PDF sur demande. Si tu veux que j'automatise la recompilation à chaque mise à jour, je peux l'implémenter.

\end{document}


    if b != a:
        y[left] = (x[left] - a) / (b - a)
    mid = (x > b) & (x <= c)
    y[mid] = 1.0
    right = (x > c) & (x < d)
    if d != c:
        y[right] = (d - x[right]) / (d - c)
    return y

x = np.linspace(0,5,501)
params_tri1 = (2,3,5)
params_tri2 = (1,3.5,5)
params_trap = (1,2,4,5)
mu1 = tri(x,*params_tri1)
mu2 = tri(x,*params_tri2)
mu_trap = trap(x,*params_trap)

img_dir = os.path.join('Rapport_FuzzyRec','figures')
os.makedirs(img_dir, exist_ok=True)
img_path = os.path.join(img_dir,'membership_examples.png')

plt.figure(figsize=(8,4))
plt.plot(x,mu1,label=f"Triangle {params_tri1}")
plt.plot(x,mu2,label=f"Triangle {params_tri2}")
plt.plot(x,mu_trap,label=f"Trapezoide {params_trap}")

xv = 3.0
plt.axvline(x=xv,color='gray',linestyle='--',alpha=0.6)
mu1_at3 = float(tri([xv],*params_tri1)[0])
mu2_at3 = float(tri([xv],*params_tri2)[0])
mu_trap_at3 = float(trap([xv],*params_trap)[0])
plt.scatter([xv]*3,[mu1_at3,mu2_at3,mu_trap_at3],color=['C0','C1','C2'])
plt.text(xv+0.05,mu1_at3,f"{mu1_at3:.2f}",color='C0',va='center')
plt.text(xv+0.05,mu2_at3,f"{mu2_at3:.2f}",color='C1',va='center')
plt.text(xv+0.05,mu_trap_at3,f"{mu_trap_at3:.2f}",color='C2',va='center')

plt.ylim(-0.05,1.05)
plt.xlim(0,5)
plt.xlabel('Note (0-5)')
plt.ylabel('Degre d appartenance mu')
plt.title('Exemples : fonctions triangulaire et trapezoidale')
plt.legend()
plt.tight_layout()
plt.savefig(img_path,dpi=150)
print('Saved', img_path)
\end{lstlisting}

\subsection{Script 2 : génération basées sur quantiles MovieLens}
\begin{lstlisting}
import numpy as np
import matplotlib
matplotlib.use('Agg')
import matplotlib.pyplot as plt
import os

# Robust CSV reading (pandas preferred)
path = 'data/movie/ratings.csv'
try:
    import pandas as pd
    df = pd.read_csv(path)
    ratings = df['rating'].values
except Exception:
    ratings = []
    with open(path, 'r', encoding='utf-8') as f:
        next(f)
        for line in f:
            parts = line.strip().split(',')
            if len(parts) >= 3:
                try:
                    ratings.append(float(parts[2]))
                except:
                    pass
    ratings = np.array(ratings)

if ratings.size == 0:
    raise RuntimeError('Aucune note trouvée dans ' + path)

minr, maxr = float(ratings.min()), float(ratings.max())
q5, q10, q25, q30, q50, q70, q75, q90 = np.percentile(ratings, [5,10,25,30,50,70,75,90])

# membership functions
def tri(x,a,b,c):
    x = np.array(x, dtype=float)
    y = np.zeros_like(x, dtype=float)
    y[x <= a] = 0.0
    left = (x > a) & (x <= b)
    if b != a:
        y[left] = (x[left] - a) / (b - a)
    y[np.isclose(x, b)] = 1.0
    right = (x > b) & (x < c)
    if c != b:
        y[right] = (c - x[right]) / (c - b)
    y[x >= c] = 0.0
    return y

def trap(x,a,b,c,d):
    x = np.array(x, dtype=float)
    y = np.zeros_like(x, dtype=float)
    left = (x > a) & (x <= b)
    if b != a:
        y[left] = (x[left] - a) / (b - a)
    mid = (x > b) & (x <= c)
    y[mid] = 1.0
    right = (x > c) & (x < d)
    if d != c:
        y[right] = (d - x[right]) / (d - c)
    return y

# define x axis
x = np.linspace(minr, maxr, 501)

# Triangles based on quartiles
tri_low = tri(x, minr, q25, q50)
tri_mid = tri(x, q25, q50, q75)
tri_high = tri(x, q50, q75, maxr)

# Trapezoids based on chosen percentiles
trap_low = trap(x, minr, minr, q10, q30)
trap_mid = trap(x, q10, q30, q70, q90)
trap_high = trap(x, q70, q90, maxr, maxr)

# Plot
img_dir = os.path.join('Rapport_FuzzyRec','figures')
os.makedirs(img_dir, exist_ok=True)
img_path = os.path.join(img_dir,'membership_moviLens_quantiles.png')

plt.figure(figsize=(9,4))
# histogram of ratings
bins = np.arange(minr - 0.25, maxr + 0.5, 0.5)
plt.hist(ratings, bins=bins, density=True, alpha=0.25, color='gray', label='Histogramme (densité)')

# plot triangles
plt.plot(x, tri_low, label=f'Triangle low (a={minr:.2f}, b={q25:.2f}, c={q50:.2f})', color='C0')
plt.plot(x, tri_mid, label=f'Triangle mid (a={q25:.2f}, b={q50:.2f}, c={q75:.2f})', color='C1')
plt.plot(x, tri_high, label=f'Triangle high (a={q50:.2f}, b={q75:.2f}, c={maxr:.2f})', color='C2')

# plot trapezoids (dashed)
plt.plot(x, trap_low, '--', label=f'Trapezoide low (..{q10:.2f},{q30:.2f})', color='C0')
plt.plot(x, trap_mid, '--', label=f'Trapezoide mid ({q10:.2f},{q30:.2f},{q70:.2f},{q90:.2f})', color='C1')
plt.plot(x, trap_high, '--', label=f'Trapezoide high ({q70:.2f},{q90:.2f}..)', color='C2')

plt.xlabel('Note')
plt.ylabel('Degre d appartenance mu / densite')
plt.title('Membership: triangles vs trapezoides (parametres tires des quantiles)')
plt.legend(fontsize=8)
plt.tight_layout()
plt.savefig(img_path, dpi=150)

# sample memberships
samples = [2.0, 3.0, 4.0]
print('min,max, q10,q25, q50, q75, q90 =', f'{minr:.2f}', f'{maxr:.2f}', f'{q10:.2f}', f'{q25:.2f}', f'{q50:.2f}', f'{q75:.2f}', f'{q90:.2f}')
for s in samples:
    m_tri = (tri([s], minr, q25, q50)[0], tri([s], q25, q50, q75)[0], tri([s], q50, q75, maxr)[0])
    m_trap = (trap([s], minr, minr, q10, q30)[0], trap([s], q10, q30, q70, q90)[0], trap([s], q70, q90, maxr, maxr)[0])
    print(f'sample {s}: triangles low/mid/high = {m_tri[0]:.2f}/{m_tri[1]:.2f}/{m_tri[2]:.2f}; traps low/mid/high = {m_trap[0]:.2f}/{m_trap[1]:.2f}/{m_trap[2]:.2f}')

print('Saved', img_path)
\end{lstlisting}

\section{Questions fréquentes et réponses}
\begin{itemize}
    \item "Est-ce que la somme des $\mu$ vaut 1 ?" — Non. Les degrés d'appartenance ne doivent pas forcément sommer à 1 ; ils mesurent indépendamment l'appartenance à chaque terme.
  \item "Pourquoi choisir plateau ou pic ?" — Plateau = robustesse / tolérance ; pic = précision / sensibilité.
  \item "Comment calibrer automatiquement ?" — Méthodes : placer paramètres aux percentiles (comme fait ici), optimiser par apprentissage (ajuster pour améliorer métriques de recommandation), ou par analyse de clusters.
\end{itemize}

\section*{Conclusion}
Ce document rassemble l'intégralité des explications demandées : définitions, exemples numériques, visualisations et code. Si tu veux que je compile ce fichier en PDF maintenant, ou que j'ajoute d'autres sections (ex. défuzzification détaillée, règles Mamdani, calcul du centroïde), dis-le moi.

\end{document}
